Scaling laws promised cost-effective experiments; six years later, they have yet to fully deliver. Instead, researchers have found them unreliable at small scales, starting at 4M parameters, and concluded that sizable models cannot be avoided. We show this is not the case and that the confounding factor is hyperparameter tuning. In fact, small models are highly sensitive, but hyperparameter sensitivity fades with scale. This small-scale sensitivity makes scaling laws easy to miss because they only emerge on the fully tuned frontier, and reaching that frontier requires an extensive search far beyond what most ever run. By ablating the basic scaling law recipe, we show well-tuned hyperparameters matter more than any other ingredient. Further, we show that as scale increases, the hyperparameter loss surface becomes lower dimensional, making it easier to find those hyperparameters. Nevertheless while scaling laws exist in small models, extrapolation hits statistical limitations and a holistic approach becomes necessary. We apply our insights to a question that once took the field years to settle: where to place normalization layers in the transformer architecture. Using a rigorous experimental setup inspired by our findings, we find that pre-normalization works better as models grow in size. With the right tools and a better understanding, small-scale experiments can deliver on scaling laws' long-awaited promise.